\section{Limitations}

The present results remain bounded in several important ways.

\begin{itemize}
\item \textbf{Paired-scanner representation learning.} SCORPION contains 48 original human slides, and the independent labeled paired-scanner cohort is canine rather than human. The experiments use frozen feature representations rather than end-to-end pixel-space adaptation. Linear probe accuracy measures recoverability under a specified estimator, not information-theoretic independence, and same-region retrieval is not a diagnostic endpoint.
\item \textbf{Whole-slide learning.} TransnnMIL has a full-PANDA stabilization result but not yet a fully matched architecture-superiority study against every relevant MIL control under identical tuning and compute. Optional hierarchical, topology, and pruning branches should not be interpreted as validated improvements until separately tested.
\item \textbf{Spatial dynamics.} The current WSI-NCA PANDA-300 experiment does not support a recurrence- or topology-dependent predictive advantage. Spatial recurrence remains an experimental extension rather than an established benefit.
\item \textbf{Federated learning.} PANDA sites are simulated partitions over pathology-derived feature data, not independent hospitals. The dominance-aware detector is validated under controlled corruption and ordinal-shift mechanisms; real institutional deployment would require prospective multi-center testing.
\item \textbf{CAMELYON17.} The source-weighting studies are centralized feature-level proxies for institutional influence and use one held-out test center. They do not validate a complete federated algorithm or establish a general law across hospitals.
\item \textbf{PCam and clinical interpretation.} PCam is a patch-level single-split benchmark, and none of the reported studies constitutes clinical validation, deployment guidance, or evidence of improved patient outcomes.
\end{itemize}

The pipeline is therefore best interpreted as a research system with several strong controlled results and several active extensions, rather than as a clinically validated end-to-end product.
